% ppo_tutorial.tex
% A tutorial on the PPO algorithm, culminating in the implementation details of PufferLib.

\documentclass[12pt, a4paper]{article}

\usepackage{amsmath, amssymb, amsthm}
\usepackage{mathtools}
\usepackage{geometry}
\usepackage{hyperref}
\usepackage{listings}
\usepackage{xcolor}
\usepackage{booktabs}
\usepackage{bm}
\usepackage{bbm}

\geometry{margin=2.5cm}

\hypersetup{
  colorlinks=true,
  linkcolor=blue!70!black,
  urlcolor=blue!70!black,
  citecolor=blue!70!black,
}

% Code listing style
\lstdefinestyle{pythonstyle}{
  language=Python,
  basicstyle=\ttfamily\small,
  keywordstyle=\color{blue!80!black}\bfseries,
  commentstyle=\color{gray},
  stringstyle=\color{orange!80!black},
  numbers=left,
  numberstyle=\tiny\color{gray},
  numbersep=8pt,
  breaklines=true,
  frame=single,
  framesep=4pt,
  backgroundcolor=\color{gray!8},
  xleftmargin=1.5em,
  framexleftmargin=1.5em,
}

\lstdefinestyle{cstyle}{
  language=C,
  basicstyle=\ttfamily\small,
  keywordstyle=\color{blue!80!black}\bfseries,
  commentstyle=\color{gray},
  stringstyle=\color{orange!80!black},
  numbers=left,
  numberstyle=\tiny\color{gray},
  numbersep=8pt,
  breaklines=true,
  frame=single,
  framesep=4pt,
  backgroundcolor=\color{gray!8},
  xleftmargin=1.5em,
  framexleftmargin=1.5em,
}

\theoremstyle{definition}
\newtheorem{definition}{Definition}[section]

\theoremstyle{plain}
\newtheorem{theorem}{Theorem}[section]
\newtheorem{lemma}[theorem]{Lemma}

\theoremstyle{remark}
\newtheorem{remark}{Remark}[section]

% Convenience macros
\newcommand{\E}{\mathbb{E}}
\newcommand{\R}{\mathbb{R}}
\newcommand{\N}{\mathbb{N}}
\newcommand{\traj}{\tau}
\newcommand{\pith}{\pi_\theta}
\newcommand{\piold}{\pi_{\theta_\mathrm{old}}}
\newcommand{\Rtraj}{R(\traj)}
\newcommand{\Jth}{J(\theta)}
\newcommand{\grad}{\nabla_\theta}
\newcommand{\KL}{\mathrm{KL}}
\newcommand{\clip}{\mathrm{clip}}

\title{\textbf{A Tutorial on Proximal Policy Optimization}\\[0.5em]
\large From Policy Gradients to PufferLib}
\author{}
\date{}

\begin{document}

\maketitle
\tableofcontents
\newpage

% ---------------------------------------------------------------------------
\section{Nomenclature and Problem Setup}
\label{sec:nomenclature}
% ---------------------------------------------------------------------------

We consider the standard reinforcement learning (RL) setting: an agent interacts with an
environment over discrete time steps, observing states and selecting actions, with the
goal of maximizing accumulated reward.

\subsection{Trajectories}

A \emph{trajectory} (or \emph{rollout}) of horizon $H$ is the sequence
\begin{equation}
  \traj = (s_0, a_0, r_0,\; s_1, a_1, r_1,\; \ldots,\; s_H, a_H, r_H),
\end{equation}
where $s_t \in \mathcal{S}$ is the state, $a_t \in \mathcal{A}$ the action, and
$r_t \in \R$ the scalar reward received at step $t$.

\subsection{Policy}

A \emph{stochastic policy} $\pith$ is a conditional distribution over actions given states,
parameterised by a vector $\theta$:
\begin{equation}
  \pith(a \mid s) = \Pr_\theta[A_t = a \mid S_t = s].
\end{equation}
In deep RL this is typically a neural network with weights $\theta$.

\subsection{Total Discounted Return}

Given a discount factor $\gamma \in [0, 1)$, the \emph{discounted return} of a
trajectory is
\begin{equation}
  \Rtraj = \sum_{t=0}^{H} \gamma^t r_t.
\end{equation}
The discount factor down-weights future rewards, both for mathematical convenience
(ensuring convergence) and to reflect temporal preference.

\subsection{Value and Advantage Functions}

\begin{definition}[State-value function]
  $V^\pi(s) = \E_\pi\!\left[\Rtraj \mid s_0 = s\right]$,
  the expected return when starting in state $s$ and following $\pi$ thereafter.
\end{definition}

\begin{definition}[Action-value function]
  $Q^\pi(s, a) = \E_\pi\!\left[\Rtraj \mid s_0 = s, a_0 = a\right]$.
\end{definition}

\begin{definition}[Advantage function]
  \begin{equation}
    A^\pi(s, a) = Q^\pi(s, a) - V^\pi(s).
  \end{equation}
  The advantage measures how much better action $a$ is compared to the average action
  taken by $\pi$ in state $s$.
\end{definition}

\subsection{Utility Function}

The quantity we want to maximise is the \emph{utility}
\begin{equation}
  \Jth = \E_{\traj \sim \pith}\!\left[\Rtraj\right]
       = \E_{\traj \sim \pith}\!\left[\sum_{t=0}^{H} \gamma^t r_t\right].
\end{equation}
Since the expectation is taken over trajectories sampled from the current policy, both
the sampling distribution \emph{and} the rewards depend on $\theta$, making
direct optimisation non-trivial.

% ---------------------------------------------------------------------------
\section{The Policy Gradient Theorem}
\label{sec:pg}
% ---------------------------------------------------------------------------

\subsection{Deriving the Gradient of the Utility}

To apply first-order (gradient) methods we need $\grad \Jth$.  The key observation is the
\emph{log-derivative trick} (also called the REINFORCE trick):

\begin{equation}
  \grad \pith(a \mid s)
    = \pith(a \mid s)\, \grad \log \pith(a \mid s).
\end{equation}

The probability of a trajectory under policy $\pith$ factorises as
\begin{equation}
  p_\theta(\traj) = p(s_0) \prod_{t=0}^{H} \pith(a_t \mid s_t)\, p(s_{t+1} \mid s_t, a_t),
\end{equation}
so the log-probability decomposes as
\begin{equation}
  \log p_\theta(\traj) = \log p(s_0)
    + \sum_{t=0}^{H} \log \pith(a_t \mid s_t)
    + \sum_{t=0}^{H} \log p(s_{t+1} \mid s_t, a_t).
\end{equation}
Only the middle sum depends on $\theta$, hence
\begin{equation}
  \grad \log p_\theta(\traj) = \sum_{t=0}^{H} \grad \log \pith(a_t \mid s_t).
\end{equation}

Applying the log-derivative trick to the utility gives
\begin{align}
  \grad \Jth
    &= \grad \int p_\theta(\traj)\, \Rtraj\, d\traj \nonumber\\
    &= \int \grad p_\theta(\traj)\, \Rtraj\, d\traj \nonumber\\
    &= \int p_\theta(\traj)\, \grad \log p_\theta(\traj)\, \Rtraj\, d\traj \nonumber\\
    &= \E_{\traj \sim \pith}\!\left[\Rtraj\, \grad \log p_\theta(\traj)\right].
\end{align}

Substituting the factorisation of $\grad \log p_\theta(\traj)$ yields the
\textbf{REINFORCE policy gradient estimator}:
\begin{equation}
  \boxed{
    \grad \Jth
      = \E_{\traj \sim \pith}\!\left[
          \sum_{t=0}^{H} \grad \log \pith(a_t \mid s_t)\, \Rtraj
        \right].
  }
  \label{eq:reinforce}
\end{equation}

\subsection{Monte-Carlo Estimation}

In practice we approximate Eq.~\eqref{eq:reinforce} with a sample average over $N$
trajectories $\{\traj^{(i)}\}_{i=1}^N$:
\begin{equation}
  \grad \Jth \approx
    \frac{1}{N}\sum_{i=1}^{N} \sum_{t=0}^{H}
      \grad \log \pith(a_t^{(i)} \mid s_t^{(i)})\, R(\traj^{(i)}).
\end{equation}
This estimator is unbiased but can have very high variance, which motivates the
modifications described in Section~\ref{sec:variance}.

% ---------------------------------------------------------------------------
\section{Variance Reduction}
\label{sec:variance}
% ---------------------------------------------------------------------------

High variance in the policy gradient estimate leads to unstable training.  Three
standard techniques are used to reduce it.

\subsection{Reward-to-Go}

Using the full return $R(\tau)$ introduces unnecessary noise: actions at step $t$
cannot influence rewards received \emph{before} $t$.  Replacing $\Rtraj$ with
the \emph{reward-to-go}
\begin{equation}
  \hat{R}_t = \sum_{t'=t}^{H} \gamma^{t'-t} r_{t'}
\end{equation}
gives an estimator with the same expectation but lower variance:
\begin{equation}
  \grad \Jth
    = \E_{\traj \sim \pith}\!\left[
        \sum_{t=0}^{H} \grad \log \pith(a_t \mid s_t)\, \hat{R}_t
      \right].
\end{equation}

\subsection{Baseline Subtraction}

Adding any function $b(s_t)$ that does not depend on the action does not bias the
gradient (since $\E_a[\nabla_\theta \log \pi_\theta(a|s)\, b(s)] = 0$), but can
reduce variance.  A natural choice is the state-value function $V^\pi(s_t)$:
\begin{equation}
  \grad \Jth
    = \E\!\left[
        \sum_{t} \grad \log \pith(a_t \mid s_t)\,
        \bigl(\hat{R}_t - V^\pi(s_t)\bigr)
      \right].
\end{equation}
The term $\hat{R}_t - V^\pi(s_t)$ approximates the advantage $A^\pi(s_t, a_t)$.
In practice we use a learned critic $V_\phi(s)$ to estimate the baseline.

\subsection{Actor-Critic Methods}
\label{sec:actor_critic}

Baseline subtraction (Section~\ref{sec:variance}) still uses a Monte-Carlo estimate
of the return: the gradient update is computed only \emph{after} a full trajectory
has been collected, and the baseline $V_\phi(s_t)$ reduces variance without
contributing to the return estimate itself.

An \emph{actor-critic} method goes one step further by letting the critic
\emph{replace} the Monte-Carlo return with a bootstrapped estimate.  The policy
(the \emph{actor}) is updated using a one-step TD advantage:
\begin{equation}
  \hat{A}_t^{\mathrm{TD}} = r_{t+1} + \gamma V_\phi(s_{t+1})(1 - d_{t+1}) - V_\phi(s_t),
  \label{eq:td_adv}
\end{equation}
where $V_\phi$ is the learned critic and $d_{t+1} \in \{0, 1\}$ is a binary flag
indicating whether $s_{t+1}$ is a terminal state (i.e.\ the episode ended after
taking action $a_t$).  When $d_{t+1} = 1$ the bootstrapped term
$\gamma V_\phi(s_{t+1})$ is zeroed out, because there is no future return beyond
a terminal state.  Both components are trained concurrently:
the actor maximises expected return via the policy gradient, while the critic
minimises a regression loss against TD targets.

\paragraph{Key difference from REINFORCE with a baseline.}
In REINFORCE with a baseline the value function $V_\phi$ appears \emph{only} as a
control variate---it does not enter the return estimate.  The policy gradient still
uses the full sum of future rewards $\hat{R}_t$, which is unbiased but high-variance.
An actor-critic method substitutes $\hat{R}_t$ with the (lower-variance but biased)
bootstrapped target in Eq.~\eqref{eq:td_adv}.  The trade-off is:
\begin{center}
\begin{tabular}{lll}
  \toprule
  & \textbf{REINFORCE + baseline} & \textbf{Actor-critic (TD)} \\
  \midrule
  Bias    & None (unbiased gradient)  & Introduced by bootstrapping \\
  Variance & High (full MC return)    & Lower (single-step TD)      \\
  Updates & After full episode        & After each step (online)    \\
  \bottomrule
\end{tabular}
\end{center}

Because the critic bootstraps from its own (imperfect) predictions, approximation
bias accumulates when the value function is inaccurate.  The choice of how many
steps to unroll before bootstrapping controls the bias-variance trade-off, which
Generalised Advantage Estimation formalises as a continuous interpolation via the
$\lambda$ parameter.

\subsection{Generalised Advantage Estimation (GAE)}
\label{sec:gae}

Schulman et al.\ \cite{schulman2016gae} introduced GAE as a low-variance, low-bias estimate of the
advantage.  For a $\lambda$-weighted mixture of $k$-step TD errors
$\delta_t^{(k)} = r_t + \gamma V(s_{t+1}) - V(s_t)$,
the GAE estimator is
\begin{equation}
  \hat{A}_t^{\mathrm{GAE}(\gamma, \lambda)}
    = \sum_{l=0}^{\infty} (\gamma\lambda)^l \delta_{t+l},
  \label{eq:gae}
\end{equation}
where $\delta_t = r_{t+1} + \gamma V(s_{t+1})(1 - d_{t+1}) - V(s_t)$ is the
one-step TD residual.  The flag $d_{t+1} \in \{0, 1\}$ is 1 if the episode
terminated at step $t{+}1$ and 0 otherwise; the factor $(1 - d_{t+1})$ zeroes
out the bootstrapped value $\gamma V(s_{t+1})$ at episode boundaries so that
returns do not bleed across episodes.

Setting $\lambda = 0$ recovers the TD(0) advantage (high bias, low variance);
$\lambda = 1$ recovers the Monte-Carlo advantage (low bias, high variance).
Intermediate values trade off between the two.

% ---------------------------------------------------------------------------
\section{Trust Region Policy Optimisation (TRPO)}
\label{sec:trpo}
% ---------------------------------------------------------------------------

\subsection{Motivation}

A fundamental problem with vanilla policy gradient is that a step too large in parameter
space can collapse performance catastrophically and irreversibly, because the new
sampling distribution is far from the one used to collect data.

\subsection{The Surrogate Objective}

Schulman et al.\ \cite{schulman2015trpo} construct a \emph{surrogate objective} that can be optimised
with the \emph{old} samples.  Define the importance-sampling ratio
\begin{equation}
  r_t(\theta) = \frac{\pith(a_t \mid s_t)}{\piold(a_t \mid s_t)}.
\end{equation}
The surrogate objective is
\begin{equation}
  L^{\mathrm{CPI}}(\theta)
    = \E_t\!\left[r_t(\theta)\, \hat{A}_t\right],
\end{equation}
which matches $\Jth$ to first order around $\theta_\mathrm{old}$ but can be
optimised with fixed samples.

\subsection{The KL Constraint}

To prevent the new policy from straying too far from the old one, TRPO solves
\begin{equation}
  \max_\theta\; L^{\mathrm{CPI}}(\theta)
  \quad \text{subject to} \quad
  \E_t\!\left[\KL\!\left(\piold(\cdot \mid s_t) \;\|\; \pith(\cdot \mid s_t)\right)\right]
    \leq \delta,
  \label{eq:trpo}
\end{equation}
where $\delta$ is a small trust-region radius.

The constraint is enforced via a conjugate-gradient solve for the natural gradient
direction followed by a line-search.  This guarantees monotonic policy improvement but
comes with significant implementation complexity:
\begin{itemize}
  \item Requires second-order information (Fisher-vector products).
  \item Conjugate gradient and back-tracking line-search add engineering overhead.
  \item Difficult to apply to architectures with shared parameters between actor and
        critic.
  \item Not directly compatible with small-batch stochastic gradient descent.
\end{itemize}
These drawbacks motivate the simpler PPO family of algorithms.

% ---------------------------------------------------------------------------
\section{From TRPO to PPO: Internalising the Constraint}
\label{sec:trpo_to_ppo}
% ---------------------------------------------------------------------------

\subsection{Lagrangian Relaxation of TRPO}

The constrained problem in Eq.~\eqref{eq:trpo} can be converted to an unconstrained
one by introducing a Lagrange multiplier $\beta \geq 0$:
\begin{equation}
  \max_\theta\;
    L^{\mathrm{CPI}}(\theta)
    - \beta\;
    \E_t\!\left[\KL\!\left(\piold(\cdot \mid s_t) \;\|\; \pith(\cdot \mid s_t)\right)\right].
  \label{eq:kl_penalty}
\end{equation}
This \emph{KL-penalty} formulation is an equivalent characterisation of TRPO: for the
right value of $\beta$ the unconstrained solution coincides with the constrained one.

The trouble is that $\beta$ must be tuned---or adapted---to match the desired
trust-region radius $\delta$, and the appropriate value varies across environments and
over the course of training.  Adaptive schemes (increasing $\beta$ when the KL exceeds
$\delta$ and decreasing it otherwise) work in practice but introduce an additional
hyper-parameter loop and make the algorithm more brittle.

More fundamentally, the KL penalty is a somewhat indirect handle on what we actually
care about.  This motivates asking: \emph{what does the KL constraint protect against,
and can we express that protection more directly in the objective?}

\subsection{What Does the KL Constraint Actually Do?}

The surrogate $L^{\mathrm{CPI}}$ relies on importance sampling: data collected under
$\piold$ are reused to estimate the gradient for $\pith$.  This is valid as long as
the two distributions are close, because the importance-sampling ratio
\begin{equation}
  r_t(\theta) = \frac{\pith(a_t \mid s_t)}{\piold(a_t \mid s_t)}
\end{equation}
remains well-behaved.  When $r_t \gg 1$ or $r_t \approx 0$, a single sampled
transition dominates or vanishes in the estimate, and the surrogate diverges from the
true objective.

The KL divergence is one measure of how far $\pith$ has moved from $\piold$, but it
acts on the full distribution.  At the level of individual action probabilities, what
matters is simply that $r_t$ stays close to 1.  In fact, to first order
\begin{equation}
  \KL\!\left(\piold(\cdot \mid s) \;\|\; \pith(\cdot \mid s)\right)
  \;\approx\;
  \tfrac{1}{2}\,\E_{a \sim \piold}\!\left[(r_t(\theta) - 1)^2\right],
\end{equation}
so the KL constraint is essentially a constraint on how much $r_t$ can deviate from 1.

This suggests a simpler route: instead of penalising the KL divergence globally,
\emph{directly prevent $r_t$ from leaving a symmetric interval
$[1 - \varepsilon,\, 1 + \varepsilon]$} around 1.

\subsection{Deriving the Clipped Objective}

Suppose we want to maximise $L^{\mathrm{CPI}}(\theta) = \E_t[r_t(\theta)\,\hat{A}_t]$
subject to $r_t(\theta) \in [1-\varepsilon, 1+\varepsilon]$.  The question is what
happens to the gradient when $r_t$ reaches the boundary of this interval.

We want a \emph{pessimistic} objective: one that gives no gradient incentive to push
$r_t$ further outside the interval than it already is.  Consider the two cases
separately.

\paragraph{Case 1: $\hat{A}_t > 0$ (the action was better than average).}
We want to increase $r_t$---but only up to $1 + \varepsilon$.  Beyond that point the
update is too large to trust.  An objective that captures this is
\begin{equation}
  \min\!\bigl(r_t\,\hat{A}_t,\; (1+\varepsilon)\,\hat{A}_t\bigr)
  = \min\!\bigl(r_t,\; 1+\varepsilon\bigr)\,\hat{A}_t.
\end{equation}
The gradient is positive (encouraging increasing $r_t$) when $r_t < 1+\varepsilon$,
and zero beyond it.

\paragraph{Case 2: $\hat{A}_t < 0$ (the action was worse than average).}
We want to decrease $r_t$---but only down to $1 - \varepsilon$.  An objective that
captures this is
\begin{equation}
  \min\!\bigl(r_t\,\hat{A}_t,\; (1-\varepsilon)\,\hat{A}_t\bigr)
  = \max\!\bigl(r_t,\; 1-\varepsilon\bigr)\,\hat{A}_t.
\end{equation}
Because $\hat{A}_t < 0$, the $\min$ over two negative quantities takes the one
closest to zero, i.e.\ it stops the ratio from being pushed below $1 - \varepsilon$.

\paragraph{Unified form.}
Both cases are captured simultaneously by
\begin{equation}
  \min\!\bigl(r_t(\theta)\,\hat{A}_t,\;
              \clip\!\bigl(r_t(\theta), 1-\varepsilon, 1+\varepsilon\bigr)\,\hat{A}_t\bigr),
\end{equation}
which equals $r_t\,\hat{A}_t$ inside the interval and flattens at the boundary
regardless of the sign of $\hat{A}_t$.  Taking the expectation over time steps yields
the \emph{clipped surrogate objective}
\begin{equation}
  L^{\mathrm{CLIP}}(\theta)
    = \E_t\!\left[
        \min\!\bigl(
          r_t(\theta)\,\hat{A}_t,\;
          \clip\!\bigl(r_t(\theta), 1-\varepsilon, 1+\varepsilon\bigr)\,\hat{A}_t
        \bigr)
      \right],
\end{equation}
which is precisely the PPO objective introduced in Section~\ref{sec:ppo}.

\begin{remark}
  $L^{\mathrm{CLIP}}$ is a lower bound on $L^{\mathrm{CPI}}$ in the region where
  $r_t \in [1-\varepsilon, 1+\varepsilon]$, and equals it there exactly.  Outside the
  interval the clipping makes the bound loose in a direction that discourages further
  deviation---the objective becomes a conservative, pessimistic surrogate rather than an
  optimistic one.  This is the key structural property that gives PPO its stability.
\end{remark}

% ---------------------------------------------------------------------------
\section{Proximal Policy Optimisation (PPO)}
\label{sec:ppo}
% ---------------------------------------------------------------------------

\subsection{The Clipped Surrogate Objective}

Schulman et al.\ \cite{schulman2017ppo} propose replacing the hard KL constraint with a soft,
\emph{clipped} surrogate objective that implicitly penalises large policy changes:
\begin{equation}
  L^{\mathrm{CLIP}}(\theta)
    = \E_t\!\left[
        \min\!\bigl(
          r_t(\theta)\,\hat{A}_t,\;
          \clip\!\bigl(r_t(\theta), 1-\varepsilon, 1+\varepsilon\bigr)\,\hat{A}_t
        \bigr)
      \right],
  \label{eq:clip}
\end{equation}
where $\varepsilon > 0$ (typically $0.1$--$0.2$) is the clipping coefficient.

The intuition is straightforward:
\begin{itemize}
  \item When $\hat{A}_t > 0$, the gradient pushes $r_t$ upward.  The clip prevents the
        ratio from exceeding $1 + \varepsilon$, discouraging an excessively large update
        on a single positive sample.
  \item When $\hat{A}_t < 0$, the gradient pushes $r_t$ downward.  The clip prevents the
        ratio from falling below $1 - \varepsilon$.
\end{itemize}
Taking the $\min$ ensures we never benefit from moving $r_t$ outside $[1-\varepsilon,
1+\varepsilon]$.

\subsection{Value Function Loss}

Alongside the policy, PPO trains a value function $V_\phi$ by minimising a clipped MSE
loss:
\begin{equation}
  L^{\mathrm{VF}}(\phi)
    = \E_t\!\left[
        \max\!\left(
          \bigl(V_\phi(s_t) - \hat{R}_t\bigr)^2,\;
          \bigl(\clip(V_\phi(s_t), V_{\phi_\mathrm{old}}(s_t) - \varepsilon_V,
                                   V_{\phi_\mathrm{old}}(s_t) + \varepsilon_V)
            - \hat{R}_t\bigr)^2
        \right)
      \right],
\end{equation}
where $\hat{R}_t = \hat{A}_t + V_{\phi_\mathrm{old}}(s_t)$ are the estimated returns and
$\varepsilon_V$ (often equal to $\varepsilon$) is the value clipping coefficient.

\subsection{Entropy Bonus}

An entropy bonus is added to encourage exploration:
\begin{equation}
  L^{\mathrm{ENT}}(\theta) = \E_t\!\left[H\!\left[\pith(\cdot \mid s_t)\right]\right],
\end{equation}
where $H[\cdot]$ denotes the Shannon entropy.

\subsection{Full PPO Objective}

The combined objective minimised at each update is
\begin{equation}
  \boxed{
    L^{\mathrm{PPO}}(\theta, \phi)
      = -L^{\mathrm{CLIP}}(\theta)
        + c_\mathrm{VF}\, L^{\mathrm{VF}}(\phi)
        - c_\mathrm{ENT}\, L^{\mathrm{ENT}}(\theta),
  }
  \label{eq:ppo_full}
\end{equation}
with scalar coefficients $c_\mathrm{VF}$ and $c_\mathrm{ENT}$ controlling the relative
weighting of the three terms.

\subsection{Implementation}

PPO is attractive because it requires only first-order optimisation (e.g.\ Adam),
makes multiple gradient steps over the same batch of collected data (controlled by
\texttt{update\_epochs}), and needs no conjugate gradient or second-order information.
The clipping mechanism serves as a lightweight trust-region surrogate.

% ---------------------------------------------------------------------------
\section{Extensions in PufferLib}
\label{sec:pufferlib}
% ---------------------------------------------------------------------------

PufferLib implements PPO with several extensions that improve sample efficiency and
training stability.  The main training loop lives in
\texttt{pufferlib/pufferl.py} (\texttt{PuffeRL.train()}, lines 332--476).

\subsection{V-Trace Importance Weighting}
\label{sec:vtrace}

When data are collected asynchronously, or when multiple gradient steps are taken over
the same batch, the policy that collected the data ($\piold$) may differ from the
current policy $\pith$.  V-Trace \cite{espeholt2018impala} corrects for this
off-policyness by incorporating importance weights into the advantage computation.

Define the clipped importance ratios
\begin{align}
  \rho_t &= \min\!\left(\bar{\rho},\; \frac{\pith(a_t \mid s_t)}{\piold(a_t \mid s_t)}\right),
  \label{eq:rho}\\
  c_t    &= \min\!\left(\bar{c},\;   \frac{\pith(a_t \mid s_t)}{\piold(a_t \mid s_t)}\right).
  \label{eq:c}
\end{align}
The V-Trace target value is
\begin{equation}
  v_t = V(s_t) + \sum_{k=t}^{t+n-1} \gamma^{k-t}
        \left(\prod_{i=t}^{k-1} c_i\right) \rho_k\,
        \delta_k,
  \label{eq:vtrace_target}
\end{equation}
where $\delta_k = r_{k+1} + \gamma V(s_{k+1}) - V(s_k)$ is the one-step TD error.

PufferLib unifies GAE (Eq.~\ref{eq:gae}) and V-Trace (Eq.~\ref{eq:vtrace_target}) into
a single backward pass.  Inspecting the CUDA kernel in
\texttt{pufferlib/extensions/cuda/pufferlib.cu} (lines 7--20):

\begin{lstlisting}[style=cstyle,
  caption={Unified GAE + V-Trace advantage kernel (CUDA, \texttt{pufferlib/extensions/cuda/pufferlib.cu}, lines 7--20).}]
__host__ __device__ void puff_advantage_row_cuda(
        float* values, float* rewards, float* dones,
        float* importance, float* advantages,
        float gamma, float lambda,
        float rho_clip, float c_clip, int horizon) {
    float lastpufferlam = 0;
    for (int t = horizon-2; t >= 0; t--) {
        int t_next = t + 1;
        float nextnonterminal = 1.0 - dones[t_next];
        float rho_t = fminf(importance[t], rho_clip);   // Eq. (16)
        float c_t   = fminf(importance[t], c_clip);     // Eq. (17)
        float delta = rho_t * (rewards[t_next]
                     + gamma * values[t_next] * nextnonterminal
                     - values[t]);                       // rho-scaled TD error
        lastpufferlam = delta
                      + gamma * lambda * c_t
                        * lastpufferlam * nextnonterminal;
        advantages[t] = lastpufferlam;
    }
}
\end{lstlisting}

The recursion is a backward sweep over the horizon.  The \texttt{importance} buffer
holds $\pi_\theta(a_t|s_t)/\pi_{\theta_\mathrm{old}}(a_t|s_t)$ ratios computed from the
most recent forward pass.  An identical CPU fallback lives in
\texttt{pufferlib/extensions/pufferlib.cpp} (lines 28--41).

The function is invoked in the training loop (\texttt{pufferl.py}, lines 354--356):

\begin{lstlisting}[style=pythonstyle,
  caption={Advantage computation call in \texttt{pufferlib/pufferl.py} (lines 352--356).}]
shape = self.values.shape
advantages = torch.zeros(shape, device=device)
advantages = compute_puff_advantage(
    self.values, self.rewards,
    self.terminals, self.ratio, advantages,
    config['gamma'], config['gae_lambda'],
    config['vtrace_rho_clip'], config['vtrace_c_clip'])
\end{lstlisting}

Relevant configuration values (from \texttt{pufferlib/config/default.ini}):
\begin{center}
\begin{tabular}{lll}
  \toprule
  Parameter & Default & Description \\
  \midrule
  \texttt{gamma}          & $0.995$ & Discount factor $\gamma$ \\
  \texttt{gae\_lambda}    & $0.90$  & GAE $\lambda$ \\
  \texttt{vtrace\_rho\_clip} & $1.0$ & $\bar{\rho}$ in Eq.~\eqref{eq:rho} \\
  \texttt{vtrace\_c\_clip}   & $1.0$ & $\bar{c}$    in Eq.~\eqref{eq:c}  \\
  \bottomrule
\end{tabular}
\end{center}

\begin{remark}
  With \texttt{vtrace\_rho\_clip} = \texttt{vtrace\_c\_clip} = 1 the clipping has
  little effect (ratios are at most 1 before any update), making the estimator
  essentially equivalent to standard GAE immediately after data collection.  The
  clipping becomes meaningful when multiple gradient epochs are performed over the
  same batch (\texttt{update\_epochs} $> 1$), where the policy has already moved away
  from the behaviour policy.
\end{remark}

\subsection{Priority Experience Replay}
\label{sec:per}

Rather than sampling minibatches uniformly, PufferLib assigns a probability to each
trajectory segment proportional to the magnitude of its advantage:
\begin{equation}
  p_i \propto |A_i|^\alpha + \epsilon,
  \label{eq:per}
\end{equation}
where $\alpha$ controls the degree of prioritisation and $\epsilon$ is a small floor.

To correct for the induced sampling bias, importance-sampling weights
\begin{equation}
  w_i = \left(N \cdot p_i\right)^{-\beta}
\end{equation}
are annealed from $\beta_0$ towards 1 over training.  The implementation is in
\texttt{pufferl.py} (lines 358--363):

\begin{lstlisting}[style=pythonstyle,
  caption={Priority sampling in \texttt{pufferlib/pufferl.py} (lines 358--363).}]
# Prioritize experience by advantage magnitude
adv = advantages.abs().sum(axis=1)
prio_weights = torch.nan_to_num(adv**a, 0, 0, 0)
prio_probs = (prio_weights + 1e-6) / (prio_weights.sum() + 1e-6)
idx = torch.multinomial(prio_probs, self.minibatch_segments)
mb_prio = (self.segments * prio_probs[idx, None])**-anneal_beta
\end{lstlisting}

The minibatch advantages are then reweighted before the loss computation:

\begin{lstlisting}[style=pythonstyle,
  caption={Priority-weighted advantage normalisation (\texttt{pufferl.py}, lines 407--409).}]
adv = mb_advantages
adv = mb_prio * (adv - adv.mean()) / (adv.std() + 1e-8)
\end{lstlisting}

Relevant configuration values:
\begin{center}
\begin{tabular}{lll}
  \toprule
  Parameter & Default & Description \\
  \midrule
  \texttt{prio\_alpha}  & $0.8$ & Prioritisation exponent $\alpha$ \\
  \texttt{prio\_beta0}  & $0.2$ & Initial IS correction $\beta_0$ \\
  \bottomrule
\end{tabular}
\end{center}

\subsection{Clipped Surrogate and Value Losses}
\label{sec:clip_impl}

The PPO clipped surrogate objective (Eq.~\ref{eq:clip}) is implemented directly
in \texttt{pufferl.py} (lines 411--420):

\begin{lstlisting}[style=pythonstyle,
  caption={PPO clipped policy and value losses (\texttt{pufferl.py}, lines 411--420).}]
# Policy gradient loss
pg_loss1 = -adv * ratio                                      # unclipped
pg_loss2 = -adv * torch.clamp(ratio,
                               1 - clip_coef, 1 + clip_coef) # clipped
pg_loss = torch.max(pg_loss1, pg_loss2).mean()

# Value function loss (with clipping)
newvalue = newvalue.view(mb_returns.shape)
v_clipped = mb_values + torch.clamp(
    newvalue - mb_values, -vf_clip, vf_clip)
v_loss_unclipped = (newvalue - mb_returns) ** 2
v_loss_clipped   = (v_clipped - mb_returns) ** 2
v_loss = 0.5 * torch.max(v_loss_unclipped, v_loss_clipped).mean()
\end{lstlisting}

The total loss combines all three terms (line 424):

\begin{lstlisting}[style=pythonstyle,
  caption={Combined PPO loss (\texttt{pufferl.py}, line 424).}]
loss = pg_loss + config['vf_coef']*v_loss - config['ent_coef']*entropy_loss
\end{lstlisting}

Relevant configuration values:
\begin{center}
\begin{tabular}{lll}
  \toprule
  Parameter & Default & Description \\
  \midrule
  \texttt{clip\_coef}    & $0.2$  & $\varepsilon$ in Eq.~\eqref{eq:clip} \\
  \texttt{vf\_clip\_coef}& $0.2$  & Value function clipping coefficient \\
  \texttt{vf\_coef}      & $2.0$  & $c_\mathrm{VF}$ in Eq.~\eqref{eq:ppo_full} \\
  \texttt{ent\_coef}     & $0.001$& $c_\mathrm{ENT}$ in Eq.~\eqref{eq:ppo_full} \\
  \bottomrule
\end{tabular}
\end{center}

\subsection{Gradient Clipping and Learning Rate Schedule}

After computing the loss, gradients are accumulated over \texttt{accumulate\_minibatches}
steps and then clipped before the optimiser step (\texttt{pufferl.py}, lines 442--446):

\begin{lstlisting}[style=pythonstyle,
  caption={Gradient clipping (\texttt{pufferl.py}, lines 442--446).}]
loss.backward()
if (mb + 1) % self.accumulate_minibatches == 0:
    torch.nn.utils.clip_grad_norm_(
        self.policy.parameters(), config['max_grad_norm'])
    self.optimizer.step()
    self.optimizer.zero_grad()
\end{lstlisting}

The learning rate is linearly annealed from \texttt{learning\_rate} to
\texttt{min\_lr\_ratio} $\times$ \texttt{learning\_rate} over the course of training
(controlled by a \texttt{CosineAnnealingLR} or linear scheduler, configured via
\texttt{anneal\_lr}).

\begin{center}
\begin{tabular}{lll}
  \toprule
  Parameter & Default & Description \\
  \midrule
  \texttt{max\_grad\_norm} & $1.5$   & Gradient clip norm \\
  \texttt{learning\_rate}  & $0.015$ & Initial Adam learning rate \\
  \texttt{anneal\_lr}      & True    & Enable LR annealing \\
  \texttt{min\_lr\_ratio}  & $0.0$   & Final LR as fraction of initial \\
  \bottomrule
\end{tabular}
\end{center}

% ---------------------------------------------------------------------------
\section{Training Diagnostics: Understanding the Log}
\label{sec:log}
% ---------------------------------------------------------------------------

PufferLib prints and logs a dashboard of scalar metrics after each update.  The
relevant quantities are computed in \texttt{pufferl.py} (lines 396--457) and prefixed
with \texttt{losses/} when sent to Wandb or Neptune.

\begin{center}
\renewcommand{\arraystretch}{1.35}
\begin{tabular}{p{3.4cm} p{5.2cm} p{5.8cm}}
  \toprule
  \textbf{Metric} & \textbf{Formula / source} & \textbf{Interpretation} \\
  \midrule
  \texttt{policy\_loss}
    & $\E_t[\max(\ldots)]$ (lines 412--414)
    & PPO clipped policy gradient loss.  Magnitude alone is not diagnostic; watch for
      trends or spikes.\\

  \texttt{value\_loss}
    & $0.5\,\E_t[\max(v^2_\mathrm{unc}, v^2_\mathrm{clip})]$ (lines 418--420)
    & Critic MSE.  Should decrease over training. \\

  \texttt{entropy}
    & $\E_t[H[\pi_\theta(\cdot|s_t)]]$ (line 422)
    & Policy entropy.  Too low $\Rightarrow$ premature convergence / insufficient
      exploration. \\

  \texttt{old\_approx\_kl}
    & $\E_t[-\log r_t(\theta)]$ (line 397)
    & First-order KL approximation using the old ratio.  Should be small. \\

  \texttt{approx\_kl}
    & $\E_t[(r_t - 1) - \log r_t]$ (line 398)
    & Second-order KL approximation (Jensen--Shannon style).  If consistently $\gg 0.01$,
      consider reducing learning rate or increasing clipping.  Can be used for early
      stopping. \\

  \texttt{clipfrac}
    & $\E_t[\mathbbm{1}(|r_t - 1| > \varepsilon)]$ (line 399)
    & Fraction of timesteps where clipping is active.  Ideally $\sim 0.1$--$0.3$;
      higher values indicate the policy is being updated too aggressively. \\

  \texttt{importance}
    & $\E_t[r_t(\theta)]$ (line 438)
    & Mean importance ratio.  Should stay close to 1; large deviations indicate
      significant policy drift within a batch. \\

  \texttt{explained\_variance}
    & $1 - \mathrm{Var}(\hat{R} - \hat{V})/\mathrm{Var}(\hat{R})$ (lines 453--456)
    & Fraction of return variance explained by the critic.  Near 1 is ideal; near 0 or
      negative means the critic is no better than a constant baseline. \\
  \bottomrule
\end{tabular}
\end{center}

\begin{remark}
  A healthy training run typically shows: \texttt{explained\_variance} rising from near
  zero to $> 0.9$; \texttt{entropy} slowly decreasing; \texttt{approx\_kl} staying
  below $0.02$; and \texttt{clipfrac} in the range $0.1$--$0.3$.  Divergence in
  \texttt{value\_loss} coupled with a drop in \texttt{explained\_variance} is often an
  early warning of instability.
\end{remark}

% ---------------------------------------------------------------------------
\section{Summary}
\label{sec:summary}
% ---------------------------------------------------------------------------

\begin{enumerate}
  \item The utility $J(\theta)$ is the expected discounted return under $\pi_\theta$.
  \item The policy gradient theorem (REINFORCE) expresses $\nabla_\theta J$ as the
        expectation of the score function times the return.
  \item Variance is reduced through reward-to-go, baseline subtraction, and GAE.
  \item TRPO constrains the KL divergence to prevent catastrophic policy updates but
        requires second-order optimisation.
  \item PPO replaces the hard KL constraint with a clipped surrogate objective,
        enabling first-order optimisation while retaining stability.
  \item PufferLib extends PPO with a unified GAE + V-Trace advantage kernel, priority
        experience replay, clipped value loss, gradient accumulation, and learning rate
        annealing.  All algorithmic components are concentrated in
        \texttt{pufferlib/pufferl.py} and
        \texttt{pufferlib/extensions/cuda/pufferlib.cu}.
\end{enumerate}

% ---------------------------------------------------------------------------
% References
% ---------------------------------------------------------------------------

\begin{thebibliography}{99}

\bibitem{schulman2015trpo}
Schulman, J., Levine, S., Moritz, P., Jordan, M., \& Abbeel, P. (2015).
\textit{Trust Region Policy Optimization}.
In \textit{Proceedings of ICML 2015}, pp.\ 1889--1897.

\bibitem{schulman2016gae}
Schulman, J., Moritz, P., Levine, S., Jordan, M., \& Abbeel, P. (2016).
\textit{High-Dimensional Continuous Control Using Generalized Advantage Estimation}.
In \textit{Proceedings of ICLR 2016}.

\bibitem{schulman2017ppo}
Schulman, J., Wolski, F., Dhariwal, P., Radford, A., \& Klimov, O. (2017).
\textit{Proximal Policy Optimization Algorithms}.
arXiv preprint arXiv:1707.06347.

\bibitem{espeholt2018impala}
Espeholt, L., Soyer, H., Munos, R., Simonyan, K., Mnih, V., Ward, T., \ldots \& Kavukcuoglu, K. (2018).
\textit{IMPALA: Scalable Distributed Deep-RL with Importance Weighted Actor-Learner Architectures}.
In \textit{Proceedings of ICML 2018}, pp.\ 1407--1416.

\bibitem{schaul2016per}
Schaul, T., Quan, J., Antonoglou, I., \& Silver, D. (2016).
\textit{Prioritized Experience Replay}.
In \textit{Proceedings of ICLR 2016}.

\end{thebibliography}

\end{document}
